\section{Introduction}
A useful next-day AQI forecast must describe what could have been issued
before the target day and how well that forecast agrees with measurements.
This distinction matters when an atmospheric model represents a grid cell,
while a resident or monitoring instrument represents a particular location.
Indian high-resolution pollution forecasting already has an established
research context \cite{jena2021delhi}. The present question is whether
India-specific station correction improves a genuine operational CAMS
baseline, and whether its uncertainty can support defensible warning choices.

AirGuard India fixes issuance at 20:00 IST for the next IST calendar day.
Its retrospective stream reconstructs measured targets and admissible
forecast features; its prospective stream freezes forecasts when actually
issued. A past model field without its producing run cannot establish an
operational forecast, and a deployed dashboard cannot establish predictive
skill. Both distinctions shape the evaluation.

The work provides three research components:
\begin{itemize}
\item a versioned station dataset with recorded rejection reasons, causal
history features and a shared CPCB-style AQI engine;
\item a locked comparison of raw CAMS, direct learning and residual
correction with rolling and spatial validation and censored-Severe diagnostics;
\item an uncertainty and advisory evaluation that reports coverage failures,
under-warning and over-warning, alongside immutable prospective verification.
\end{itemize}
The second and third components are evaluation commitments rather than
completed primary findings. This manuscript distinguishes existing measured-data
evidence from pending CAMS experiments and makes no world's-first claim.
